\section{Literature Survey}
Table \ref{tab:lit} summarizes nineteen relevant studies. The comparison describes their methods, not an artificial cross-dataset accuracy ranking. The current project does not claim to invent convolution, recurrence, attention or human review.
\begin{singlespace}\small
\begin{longtable}{p{0.24\textwidth}p{0.32\textwidth}p{0.34\textwidth}}
\caption{Relevant literature and its relationship to this project}\label{tab:lit}\\
\toprule Study & Focus & Project relationship\\\midrule\endfirsthead
\toprule Study & Focus & Project relationship\\\midrule\endhead
Dubey et al. \cite{dubey2025phishx} & Character CNN plus 36-feature LightGBM & Direct reference for the first independent reproduction; not the current recurrent verdict.\\
Khonji et al. \cite{khonji2013survey} & Survey of phishing detection & Frames the diversity of URL, content and user-facing approaches.\\
Ma et al. \cite{ma2009blacklists} & Learning suspicious URL patterns beyond lists & Motivates generalization beyond exact known destinations.\\
Garera et al. \cite{garera2007framework} & Detection and measurement of phishing & Establishes URL-related evidence as an earlier research direction.\\
Blum et al. \cite{blum2010lexical} & Lexical URL features and online learning & Adaptation is prior work; our reviewed workflow is a prototype, not a new learning theory.\\
Zhang et al. \cite{zhang2007cantina} & Content-based phishing detection & Content can supplement URL evidence, but our text parser does not reproduce its detector.\\
Xiang et al. \cite{xiang2011cantina} & Feature-rich content learning & Demonstrates richer evidence; our standalone HTML signals are uncalibrated.\\
Sahoo et al. \cite{sahoo2017survey} & Malicious URL detection survey & Provides a taxonomy of representation and data challenges.\\
Saxe and Berlin \cite{saxe2017expose} & Character embeddings and convolution & Related character-pattern learning without manual tokenization.\\
Le et al. \cite{le2018urlnet} & Learned character and word URL representations & Our model has no URLNet word branch; representation choices differ.\\
Bahnsen et al. \cite{bahnsen2017rnn} & Recurrent phishing URL classification & Sequence modeling for URLs predates our BiLSTM extension.\\
Sahingoz et al. \cite{sahingoz2019machine} & Machine learning from URLs & Shows the role of extracted URL representations and comparative evaluation.\\
Jain and Gupta \cite{jain2018client} & Client-side machine-learning detection & Related deployment setting; this project uses a server API.\\
Mohammad et al. \cite{mohammad2014predicting} & Self-structuring neural detection & Neural phishing classification is established prior work.\\
Fette et al. \cite{fette2007email} & Learning phishing email detection & Our email tool is not a trained equivalent and has no measured email accuracy.\\
Zhang et al. \cite{zhang2015character} & Character CNN text classification & Motivates local character-feature extraction.\\
Ke et al. \cite{ke2017lightgbm} & Efficient gradient-boosted trees & Supplies the established engineered-feature model.\\
Hochreiter and Schmidhuber \cite{lstm} & Gated long short-term memory & Basis for the implemented recurrent component.\\
Schuster and Paliwal \cite{brnn} & Bidirectional recurrent networks & Basis for using both sequence directions.\\\bottomrule
\end{longtable}\end{singlespace}

\section{Existing System and Gap}
The reference article reports a 60/40 CNN--LightGBM ensemble, approximately 99,361 URLs, and 99.819\% accuracy on its test. It describes 70 symbols and embedding width 128. The released implementation and our original reproduction differ in some metadata; our historical encoder has 49 explicit characters and embedding width 16. This report uses source code for our implementation, rather than copying the reference architecture values into our tables.

The reference identifies adversarial robustness and richer content/context as future directions. Our extensions address bounded versions of those directions. Deterministic strings provide reproducible stress tests, while pasted HTML/email provides supplementary evidence without unsafe fetching. The newest revision adds recurrence and attention but does not establish that either is a novel algorithm.

\section{Evaluation and Calibration Lessons}
Tranco is designed as a research-oriented ranking \cite{lepochat2019tranco}; a popular root is still not equivalent to every legitimate path. Neural confidence also need not be calibrated \cite{guo2017calibration}. Adding a sigmoid supplies a probability-shaped score, not proof of calibration. Learned attention weights \cite{attention} are aggregation parameters, not a causal explanation of why a website is malicious. These distinctions guide the implementation and the claims made in later chapters.